% ==============================================================================
% Plantilla Institucional de Trabajo Terminal — UPIIZ IPN
% Archivo: capitulos/06-planteamiento-problema.tex
% ==============================================================================

\chapter{Planteamiento del Problema}
\label{cap:planteamiento}

En este capítulo se define de manera formal y precisa la problemática de ingeniería que motiva el Trabajo Terminal, las restricciones operativas del sistema y la solución mecatrónica interdisciplinaria propuesta.

% ------------------------------------------------------------------------------
\section{Formulación del problema de ingeniería}
\label{sec:planteamiento-problema}

Redacte de forma clara los síntomas y causas de la problemática que se pretende resolver. Identifique las limitaciones que presentan los sistemas actuales y formule la necesidad concreta de diseño o automatización:
\begin{enumerate}
    \item \textbf{Deficiencia o limitación principal:} Describa qué fenómeno físico, ineficiencia de proceso o restricción operacional impide alcanzar el desempeño óptimo.
    \item \textbf{Causas técnicas identificadas:} Explique los factores electromecánicos, perturbaciones ambientales o carencias sensoriales asociadas.
    \item \textbf{Consecuencias de no resolver el problema:} Puntualice los riesgos de falla, pérdidas económicas o desviaciones de calidad.
\end{enumerate}

% ------------------------------------------------------------------------------
\section{Solución mecatrónica propuesta}
\label{sec:solucion-propuesta}

Exponga la propuesta técnica de solución, detallando explícitamente cómo participará cada una de las cuatro áreas de la Ingeniería Mecatrónica:
\begin{itemize}
    \item \textbf{Mecánica:} Diseño del mecanismo, análisis cinemático y dinámico, cálculo de esfuerzos estructurales y modelado CAD 3D de piezas y ensambles.
    \item \textbf{Electrónica:} Selección y acondicionamiento de sensores, diseño de etapas de potencia, aislamiento galvánico y gestión eficiente de la energía.
    \item \textbf{Control:} Formulación de modelos dinámicos, diseño de observadores/estimadores y síntesis de leyes de control en lazo cerrado para garantizar estabilidad y seguimiento.
    \item \textbf{Programación:} Implementación de firmware determinista en tiempo real, comunicación por buses industriales/digitales e interfaz de supervisión telemática.
\end{itemize}
